\documentclass[a4paper,12pt]{article}
\usepackage[T2A]{fontenc}
\usepackage[utf8]{inputenc}
\usepackage[russian]{babel}
\usepackage{amsmath,amssymb,mathtools}
\usepackage{helvet}
\renewcommand{\familydefault}{\sfdefault}
\usepackage[margin=21mm,top=22mm,bottom=24mm,headheight=14pt]{geometry}
\usepackage{microtype}
\usepackage{array,booktabs,tabularx}
\usepackage{enumitem}
\usepackage{xcolor}
\usepackage{hyperref}
\usepackage{tikz}
\usepackage{fancyhdr}
\usepackage{needspace}
\definecolor{accent}{HTML}{006F73}
\definecolor{ink}{HTML}{253034}
\definecolor{error}{HTML}{A02B30}
\definecolor{hint}{HTML}{235C84}
\definecolor{rulegray}{HTML}{BCC8C8}
\hypersetup{colorlinks=true,linkcolor=accent,pdfauthor={Лёвин Артём Александрович},pdftitle={Ревью домашней работы — 26 октября 2026}}
\setlength{\parindent}{0pt}
\setlength{\parskip}{6pt}
\setlength{\emergencystretch}{2em}
\renewcommand{\arraystretch}{1.35}
\pagestyle{fancy}
\fancyhf{}
\fancyhead[L]{\small\color{ink!75}Ревью домашней работы}
\fancyhead[R]{\small\color{ink!75}26.10.2026}
\fancyfoot[C]{\scriptsize\color{ink!70}Лёвин Артём Александрович \quad\textbullet\quad эксклюзивно для Марка Гукина \quad\textbullet\quad \thepage}
\renewcommand{\headrulewidth}{0.3pt}
\renewcommand{\footrulewidth}{0pt}
\newcommand{\reviewlabel}[1]{\par\smallskip\textbf{\color{accent}#1}\par\nobreak}
\newcommand{\errortext}{\textbf{\color{error}Ошибка.}\ }
\newcommand{\ideatext}{\textbf{\color{hint}Ключевая идея.}\ }
\newcommand{\resulttext}{\textbf{\color{accent}Итог.}\ }
\newcommand{\sep}{\vspace{4pt}\par\noindent\color{rulegray}\rule{\linewidth}{0.35pt}\color{black}\vspace{4pt}}
\newcommand{\tasktitle}[2]{\clearpage\phantomsection\label{task:#1}\section*{Задание №\,#1\,\,\textcolor{accent}{\normalfont\large #2}}\vspace{-3pt}\sep}
\begin{document}
\thispagestyle{empty}
\begin{center}
\vspace*{35mm}
{\large\color{accent}МАТЕМАТИКА}\par
\vspace{20mm}
{\fontsize{29}{37}\selectfont\bfseries Ревью домашней работы}\par
\vspace{16mm}
{\Large 26 октября 2026 г.}\par
\vspace{30mm}
{\large Лёвин Артём Александрович}\par
\end{center}
\vfill
\begin{center}
\begin{tikzpicture}[x=1cm,y=1cm]
\draw[accent!55,line width=0.9pt] (-6,0) .. controls (-3.9,0.54) and (-2.2,-0.45) .. (0,0)
 .. controls (2.2,0.45) and (3.9,-0.54) .. (6,0);
\end{tikzpicture}
\end{center}
\newpage
\section*{Сводная таблица ошибок}
\vspace{3mm}
\small
\begin{tabularx}{\linewidth}{@{}>{\raggedright\arraybackslash}p{10mm}>{\raggedright\arraybackslash}p{30mm}X>{\raggedright\arraybackslash}p{27mm}>{\raggedright\arraybackslash}p{28mm}@{}}
\toprule
\textbf{№} & \textbf{Тема} & \textbf{Суть ошибки} & \textbf{Ваш ответ} & \textbf{Правильный ответ}\\
\midrule
\hyperref[task:1]{\textbf{1}} & Степени и корни & Потерян множитель $\sqrt[3]{8}$ в знаменателе. & $8$ & $4$\\
\addlinespace[3pt]
\hyperref[task:3]{\textbf{3}} & Показательное уравнение & При переносе единицы не изменена правая часть. & $x=-\frac23$ & $x=-\frac76$\\
\addlinespace[3pt]
\hyperref[task:4]{\textbf{4}} & Дробная степень & Потерян отрицательный корень. & $x=27$ & $x=-27;\ 27$\\
\addlinespace[3pt]
\hyperref[task:10]{\textbf{10}} & Логарифмы и ОДЗ & Параметры найдены, но уравнение осталось нерешённым. & $a=\sqrt5$, $b=4$; $x$ не найден & $a=\sqrt5$, $b=4$; решений нет\\
\bottomrule
\end{tabularx}
\normalsize
\vfill
\begin{center}\small\color{ink!70}В разбор включены только задания, требующие исправления или завершения.\end{center}
\clearpage
\thispagestyle{empty}
\vspace*{0.36\textheight}
\begin{center}
{\fontsize{26}{32}\selectfont\bfseries Разбор ошибок}\par
\vspace{7mm}
{\large\color{accent}Шаг за шагом}\par
\end{center}
\clearpage

\tasktitle{1}{Степени и корни}
\reviewlabel{Текст задания}
Вычислите без калькулятора:
\[
E=\frac{8^{\sqrt3+2}\cdot4^{1-\frac32\sqrt3}}{2^5\sqrt[3]{8}}.
\]
\reviewlabel{Ваше решение}
Основания $8$ и $4$ заменены степенями двойки. Затем получено $2^8/2^5=2^3=8$.
\reviewlabel{Ваш ответ}
$8$.
\reviewlabel{Ошибка}
\errortext При упрощении знаменателя потерян множитель $\sqrt[3]{8}=2$.
\reviewlabel{Где именно возникает ошибка}
Последний правильный шаг --- приведение всех множителей к основанию $2$:
\[
8=2^3,\qquad4=2^2,\qquad\sqrt[3]{8}=2.
\]
В знаменателе исходного выражения находятся \emph{два множителя}: $2^5$ и $\sqrt[3]{8}$. В Вашем следующем преобразовании остаётся только $2^5$.
\reviewlabel{Почему это неверно}
Знаменатель равен $2^5\cdot2=2^6$. Пропуск множителя $2$ уменьшает знаменатель вдвое, вследствие чего значение всей дроби увеличивается вдвое. Это и приводит к ответу $8$ вместо $4$.
\reviewlabel{Ключевая идея}
\ideatext Преобразовывайте в степень двойки каждый множитель числителя \emph{и} знаменателя, включая корень. Только после этого складывайте и вычитайте показатели.
\reviewlabel{Исправление от последнего правильного шага}
Оставляем найденный показатель числителя $8$ и возвращаем пропущенный множитель:
\[
\frac{2^8}{2^5\cdot\boxed{2}}=\frac{2^8}{2^6}=2^2=4.
\]
\reviewlabel{Полное правильное решение}
\begin{align*}
E&=\frac{(2^3)^{\sqrt3+2}(2^2)^{1-\frac32\sqrt3}}{2^5\cdot2^{\frac33}}\\[2pt]
&=\frac{2^{3\sqrt3+6}\cdot2^{2-3\sqrt3}}{2^{5+1}}\\[2pt]
&=2^{3\sqrt3+6+2-3\sqrt3-6}\\[2pt]
&=2^2=4.
\end{align*}
\reviewlabel{Проверка}
Иррациональные слагаемые $3\sqrt3$ и $-3\sqrt3$ взаимно уничтожаются. В числителе остаётся $2^8$, в знаменателе --- $2^6$, поэтому результат равен $4$.
\reviewlabel{Итог}
\resulttext $\boxed{E=4}$.
\reviewlabel{Задача на закрепление}
Вычислите без калькулятора:
\[
\frac{8^{\sqrt2+1}\cdot4^{2-\frac32\sqrt2}}{2^3\sqrt[3]{8}}.
\]

\tasktitle{3}{Показательное уравнение}
\reviewlabel{Текст задания}
Решите уравнение:
\[
9\cdot 3^{2x-1}=\frac1{\sqrt[3]{81}}.
\]
\reviewlabel{Ваше решение}
Вы получили равенство $3^{2x+1}=3^{-4/3}$ и перешли к показателям:
\[
2x+1=-\frac43.
\]
Затем записали $2x=-\frac43$ и получили $x=-\frac23$.
\reviewlabel{Ваш ответ}
$x=-\frac23$.
\reviewlabel{Ошибка}
\errortext При переходе от $2x+1=-\frac43$ к $2x=-\frac43$ пропущено вычитание единицы из правой части.
\reviewlabel{Где именно возникает ошибка}
Последний правильный шаг:
\[
2x+1=-\frac43.
\]
Первый неправильный шаг:
\[
2x=-\frac43.
\]
\reviewlabel{Почему это неверно}
Чтобы оставить слева $2x$, необходимо вычесть $1$ \emph{из обеих частей} уравнения. Если изменить только левую часть, равенство перестаёт быть равносильным исходному.
\reviewlabel{Ключевая идея}
\ideatext При переносе слагаемого через знак равенства меняется его знак. Для дробей предварительно приводите целое число к общему знаменателю.
\reviewlabel{Исправление от последнего правильного шага}
\begin{align*}
2x+1&=-\frac43,\\
2x&=-\frac43-1=-\frac43-\frac33=-\frac73,\\
x&=-\frac76.
\end{align*}
\reviewlabel{Полное правильное решение}
Представим обе части как степени тройки:
\[
9=3^2,\qquad81=3^4,\qquad
\frac1{\sqrt[3]{81}}=(3^4)^{-1/3}=3^{-4/3}.
\]
Тогда
\begin{align*}
3^2\cdot3^{2x-1}&=3^{-4/3},\\
3^{2x+1}&=3^{-4/3},\\
2x+1&=-\frac43,\\
2x&=-\frac73,\\
x&=-\frac76.
\end{align*}
Сравнение показателей допустимо, поскольку основание $3$ положительно и не равно единице.
\reviewlabel{Проверка}
При $x=-\frac76$ имеем $2x-1=-\frac{10}{3}$, поэтому
\[
9\cdot3^{-10/3}=3^{2-10/3}=3^{-4/3}
=\frac1{\sqrt[3]{81}}.
\]
\reviewlabel{Итог}
\resulttext $\boxed{x=-\frac76}$.
\reviewlabel{Задача на закрепление}
Решите уравнение:
\[
9\cdot3^{2x+1}=\frac1{\sqrt[3]{27}}.
\]

\tasktitle{4}{Уравнение с дробной степенью}
\reviewlabel{Текст задания}
Найдите все действительные корни уравнения:
\[
x^{-2/3}=\frac19.
\]
\reviewlabel{Ваше решение}
После преобразования дробной степени Вы получили $x^{2/3}=9$ и далее $x=3^3=27$.
\reviewlabel{Ваш ответ}
$x=27$.
\reviewlabel{Ошибка}
\errortext Потерян отрицательный корень: уравнение с квадратом может иметь два действительных решения.
\reviewlabel{Где именно возникает ошибка}
Последний правильный шаг:
\[
x^{2/3}=9.
\]
Первый неполный шаг --- переход сразу к $\sqrt[3]{x}=3$, без рассмотрения значения $-3$.
\reviewlabel{Почему это неверно}
В действительных числах
\[
x^{2/3}=(\sqrt[3]{x})^2.
\]
Квадраты чисел $3$ и $-3$ одинаковы и равны $9$. Кубный корень определён также для отрицательных чисел, поэтому необходимо рассмотреть обе возможности.
\reviewlabel{Ключевая идея}
\ideatext Если $u^2=c$, где $c>0$, то $u=\pm\sqrt c$. При дробной степени с чётным числителем этот выбор знака особенно важен.
\reviewlabel{Исправление от последнего правильного шага}
\begin{align*}
x^{2/3}&=9,\\
(\sqrt[3]{x})^2&=9,\\
\sqrt[3]{x}&=3\quad\text{или}\quad\sqrt[3]{x}=-3,\\
x&=27\quad\text{или}\quad x=-27.
\end{align*}
\reviewlabel{Полное правильное решение}
Исходное выражение определено при $x\ne0$, так как показатель отрицательный:
\[
x^{-2/3}=\frac{1}{(\sqrt[3]{x})^2}.
\]
Для $x\ne0$ знаменатель положителен, и можно выполнить равносильные преобразования:
\begin{align*}
\frac1{(\sqrt[3]{x})^2}&=\frac19,\\
(\sqrt[3]{x})^2&=9,\\
\sqrt[3]{x}&=\pm3,\\
x&=(\pm3)^3=\pm27.
\end{align*}
Оба значения отличны от нуля и допустимы.
\reviewlabel{Проверка}
\[
27^{-2/3}=\frac1{3^2}=\frac19,\qquad
(-27)^{-2/3}=\frac1{(-3)^2}=\frac19.
\]
\reviewlabel{Итог}
\resulttext $\boxed{x\in\{-27,\,27\}}$.
\reviewlabel{Задача на закрепление}
Найдите все действительные корни уравнения:
\[
x^{-2/3}=\frac1{16}.
\]

\tasktitle{10}{Логарифмическая функция и ОДЗ}
\reviewlabel{Текст задания}
График функции $f(x)=\log_a(x+b)$ проходит через точки $A(-3;0)$ и $B(1;2)$. Найдите $a$, $b$, а затем решите уравнение
\[
f(x)=\log_a(2x+7)-2.
\]
Выпишите ОДЗ до преобразований.
\reviewlabel{Ваше решение}
Вы верно получили $b=4$ и $a=\sqrt5$ и указали ограничение $2x+7>0$ (то есть $x>-\frac72$). В конце записано $f(x)=\log_{\sqrt5}(x+4)$; решение последующего уравнения отсутствует.
\reviewlabel{Ваш ответ}
$a=\sqrt5$, $b=4$; ответ относительно $x$ не указан.
\reviewlabel{Ошибка}
\errortext Задание выполнено частично: параметры определены правильно, но равенство двух логарифмических выражений не решено и итоговая проверка ОДЗ не проведена.
\reviewlabel{Где именно возникает ошибка}
Последний правильный шаг:
\[
f(x)=\log_{\sqrt5}(x+4).
\]
Следующий необходимый шаг --- подставить эту формулу в заданное уравнение:
\[
\log_{\sqrt5}(x+4)=\log_{\sqrt5}(2x+7)-2.
\]
Этого перехода и дальнейшего отбора корней в записи нет.
\reviewlabel{Почему это важно}
Найти параметры функции --- только первая часть задания. При решении логарифмического уравнения нужно заранее определить область допустимых значений, выполнить равносильные преобразования и проверить, что найденное число принадлежит ОДЗ. Алгебраическое преобразование может привести к числу вне области определения.
\reviewlabel{Ключевая идея}
\ideatext Сначала выпишите условия положительности аргументов \emph{обоих} логарифмов, затем представьте свободный член $2$ как логарифм по тому же основанию и сравните аргументы.
\reviewlabel{Исправление от последнего правильного шага}
Подставляем найденные параметры и объединяем условия:
\[
\begin{cases}x+4>0,\\2x+7>0,\end{cases}
\qquad\Longleftrightarrow\qquad x>-\frac72.
\]
Поскольку $(\sqrt5)^2=5$, имеем $2=\log_{\sqrt5}5$. Тогда
\begin{align*}
\log_{\sqrt5}(x+4)&=\log_{\sqrt5}(2x+7)-\log_{\sqrt5}5,\\
\log_{\sqrt5}(x+4)&=\log_{\sqrt5}\!\left(\frac{2x+7}{5}\right),\\
x+4&=\frac{2x+7}{5},\\
5x+20&=2x+7,\\
3x&=-13,\\
x&=-\frac{13}{3}.
\end{align*}
Полученное число меньше $-\frac72$, поэтому оно недопустимо.
\reviewlabel{Полное правильное решение}
\textbf{1. Найдём параметры.} Для точки $A(-3;0)$:
\[
0=\log_a(b-3)\quad\Longrightarrow\quad b-3=1
\quad\Longrightarrow\quad b=4.
\]
Для точки $B(1;2)$:
\[
2=\log_a(1+4)=\log_a5
\quad\Longrightarrow\quad a^2=5.
\]
Основание логарифма должно быть положительным и отличным от единицы, поэтому $a=\sqrt5$.

\textbf{2. Запишем ОДЗ.} Для исходного уравнения при $a=\sqrt5$:
\[
\begin{cases}x+4>0,\\2x+7>0,\end{cases}
\quad\Longleftrightarrow\quad x>-\frac72.
\]

\textbf{3. Решим уравнение.} Используем $2=\log_{\sqrt5}5$:
\[
\log_{\sqrt5}(x+4)=\log_{\sqrt5}\!\left(\frac{2x+7}{5}\right).
\]
Из равенства логарифмов с одним допустимым основанием следует равенство аргументов:
\[
5(x+4)=2x+7
\quad\Longrightarrow\quad x=-\frac{13}{3}.
\]
\textbf{4. Проверим принадлежность ОДЗ.} Полученный кандидат удовлетворяет
\[
-\frac{13}{3}<-\frac72.
\]
Значит, аргумент $x+4$ уже отрицателен, и логарифм $\log_{\sqrt5}(x+4)$ при этом $x$ не определён в действительных числах.
\reviewlabel{Проверка}
При $x=-\frac{13}{3}$:
\[
x+4=-\frac13<0,\qquad 2x+7=-\frac53<0.
\]
Оба логарифма не определены; иных кандидатов линейное уравнение не даёт.
\reviewlabel{Итог}
\resulttext $\boxed{a=\sqrt5,\ b=4;\quad x\in\varnothing.}$
\reviewlabel{Задача на закрепление}
График функции $g(x)=\log_a(x+b)$ проходит через точки $A(-2;0)$ и $B(1;2)$. Найдите $a$ и $b$, затем решите уравнение
\[
g(x)=\log_a(3x+4)-1.
\]
Выпишите ОДЗ до преобразований.

\end{document}
